\documentclass[10pt,a4paper]{amsart}
\usepackage[margin=19mm]{geometry}
\usepackage{amsmath,amssymb,mathtools}
\usepackage[scr=rsfs,cal=euler]{mathalfa}
\usepackage{xfrac}
\usepackage[unicode,hidelinks,bookmarks=false]{hyperref}
\newcommand{\ie}{i.e.\ }
\newcommand{\cf}{cf.\ }
\newcommand{\resp}{respectively\ }
\newcommand{\Subsection}{\subsection}
\providecommand{\nobreakdash}{\nobreak\hbox{-}}
\setlength{\emergencystretch}{2em}
\multlinegap0pt
\usepackage{amssymb}
\usepackage[normalem]{ulem}
\newtheorem{lemma}{Lemma}
\newcommand{\De}{\Delta}
\newcommand{\MM}{\mathcal{M}} %for manifold of const curvature metrics
\newcommand{\N}{\ensuremath{{\mathbb N}}}
\newcommand{\R}{\ensuremath{{\mathbb R}}}
\newcommand{\Si}{\Sigma}
\newcommand{\beq}{\begin{equation}}
\newcommand{\beqas}{\begin{equation*}\begin{aligned}}
\newcommand{\beqs}{\begin{equation*}}
\newcommand{\dd}{\mathrm{d}} %d for differentials
\newcommand{\eeq}{\end{equation}}
\newcommand{\eeqas}{\end{aligned}\end{equation*}}
\newcommand{\eeqs}{\end{equation*}}
\newcommand{\ep}{\varepsilon}
\newcommand{\eps}{\varepsilon}
\newcommand{\half}{\frac{1}{2}}
\newcommand{\hf}{\frac{1}{2}} % shorthand half
\DeclareMathOperator{\inj}{inj}
\newcommand{\inner}[2]{\langle #1, #2\rangle} % inner product
\newcommand{\la}{\lambda} % shorthand lambda
\newcommand{\na}{\nabla}
\newcommand{\norm}[1]{\Vert#1\Vert} 
\newcommand{\p}{\partial} % partial symbol
\newcommand{\pSi}{{\partial \Sigma}}
\newcommand{\vph}{\varphi}
\newcommand{\vphi}{\varphi}
\title{\sc Flowing to free boundary minimal surfaces}
\author{Melanie Rupflin, Michael Struwe, Christopher Wright}
\date{Source: arXiv:2605.19574v1 (CC BY 4.0)}
\begin{document}
\maketitle

\noindent\emph{Source and licence.} \sc Flowing to free boundary minimal surfaces. Version arXiv:2605.19574v1 (CC BY 4.0), distributed by arXiv under the Creative Commons Attribution 4.0 International licence, http://creativecommons.org/licenses/by/4.0/, as recorded in the arXiv licence metadata for this exact version and shown on the abstract page https://arxiv.org/abs/2605.19574v1. Full licence notice: \texttt{licenses/arxiv-2605.19574-CC-BY-4.0.txt}.\par\smallskip

\noindent\textit{Editorial scope.} Counted unit: the article's lemma lemma:ep\_time\_ind\_weak\_exist (source0.tex lines 951--959). The article's complete original proof of it is delivered with it (source0.tex lines 961--1033, one \texttt{proof} environment of 73 lines). No mathematics is written, repaired or re-derived: the statement and the proof are the article's own lines, in the article's own notation, and no retained line is substituted or re-flowed.  Retained uncounted as the source-local context the proof needs: lines 484--486.  Nothing was omitted from any retained span, so the package carries no omission record.  No retained line contains a citation, so the wrapper carries no bibliography and no bibliography entry.  No retained line contains a figure environment, an image-inclusion command or drawing code, so no image is delivered and the package declares no asset dependency.  Editorial formatting only, in addition: an AMS \texttt{amsart} preamble that restates the article's own declarations and macros used by the retained lines, namely the environments \texttt{lemma} on separate counters, together with the standard packages those lines need.  The retained environments print in this document as Lemma 1 in order of appearance, whereas the article prints the same items with the numbering of its own full document.  The complete native source of the article as distributed in the arXiv e-print for this exact version is bundled unchanged as \texttt{sources/200924/source0.tex}.
\begin{equation}\label{eq:ep_v_map}
	\p_t u = -(\ep + P_{v})\p_{\nu_{g}}u_{g}\ \hbox{ on }\pSi
\end{equation}

\begin{lemma}\label{lemma:ep_time_ind_weak_exist}
    Let $\ep\in (0, 1/2]$, $T>0$,  $g\in\MM^{3}$ and let $v\in H^{1/2}(\pSi;N_\eta)$. Then for any $u_0 \in  H^{\hf}(\p\Si;\R^n)$ there exists a unique weak solution $u:[0,T]\times \pSi\to \R^n$ of \eqref{eq:ep_v_map} with 
    $$u\in H^1([0,T];L^2(\p\Si))\cap L^2([0,T];H^1(\p\Si)) \text{ and } u_g\in  L^\infty([0,T];H^1(\Si)) %, \text{ and } \p_{\nu} u_g\in L^2([0,T]\times\p\Si)
    $$
    whose $L^2(\pSi)$ trace at $t=0$ is given by $u_0$. Furthermore, along any weak solution with this regularity we have  
    \beq\label{eq:energy-dec_ep_approx}
        \frac{\dd}{\dd t}E_\half(u,g)\leq -\half\norm{\p_tu}_{L^2(\p\Si)}^2.
    \eeq
\end{lemma}

\begin{proof}[Proof of Lemma \ref{lemma:ep_time_ind_weak_exist}]
    As $g$ is independent of time, we can consider a fixed orthonormal basis $\vph_0,\vph_1,\ldots$ of $L^2(\pSi)$ which consists of Steklov eigenfunctions corresponding to non-decreasing eigenvalues $0 = \la_0 < \la_1 \leq \la_2 \leq \cdots $, i.e. functions $\vph_j:\pSi\to \R$ with
    \begin{align*}
    	\p_{\nu}\big[ (\vph_j)_g\big] = \la_j \vph_j \text{ on }\p\Si,
    \end{align*}
    and work with the $\R^n$-valued spans
    \beqs
        V_\ell := \big\{u:\pSi\to \R^n: u=\sum_{j=0}^\ell\varphi_j b_j \text{ for some } b_j\in\R^n\big\},\quad \ell \in \N.
    \eeqs
    For any $\ell\in\N$, we consider the functions $u_\ell(t)(x)= \sum_{j=0}^\ell a_{j,\ell}(t)\vph_j(x) \in C^\infty([0,T];V_\ell)$, and their harmonic extensions $U_\ell=(u_\ell)_g$, which are chosen so that  $u_\ell(0)$ is given as the $L^2(\pSi)$-orthonormal projection of $u_0$ onto $V_\ell$ and which are so that
    \beq\label{eq:ep_galerkin_b}
        \int_{\p\Si}\p_t u_\ell\cdot w\dd s = \int_{\p\Si}Q(\p_{\nu} U_\ell) \cdot w\dd s\ \text{ for all }w\in V_\ell
    \eeq
    holds at each time in $[0,T]$. As $\p_{\nu}U_\ell(t)\vert_{\pSi} = \sum_{j=0}^\ell a_{j,\ell}(t)\la_j \vph_j \in V_\ell$ at each such time, and as the $\vph_j$ are orthonormal, we can obtain these functions $u_\ell$ simply by solving the corresponding system of ODEs %ordinary differential equations
    \beqs
        \p_ta_{j,\ell} = \sum_{i=0}^\ell \la_i\int_{\p\Si}Q(a_{i,\ell})\vph_i\vph_j\dd s,\ 0\le j\le\ell,
    \eeqs
    %for the coefficients $a_{j,\ell}$ 
    with the initial condition
    \beqs
        a_{j,\ell}(0) =a_j(u_0):= \int_{\p\Si}u_0\vph_j \dd s\in \R^n,\ 0\le j\le\ell.
    \eeqs
    As \eqref{eq:ep_galerkin_b} is in particular applicable for $v=\p_\nu u_\ell(t)\in V_\ell$ and as $\De U_\ell=0$ we note that the energy of these Galerkin approximates decays according to     
    \begin{align}\label{est:galerkin_uniform_bound}
        \hf\frac{\dd}{\dd t}\norm{\nabla U_\ell}_{L^2(\Si)}^2 &= \int_\Si\inner{\nabla\p_t U_\ell}{\nabla U_\ell}\dd v = \int_{\p\Si}\p_tu_\ell\cdot\p_{\nu}U_\ell\dd s\nonumber\\
        & = \int_{\p\Si}Q(\p_{\nu}U_\ell)\cdot\p_{\nu}U_\ell\dd s = -\ep\norm{\p_{\nu}U_\ell}_{L^2(\p\Si)}^2-\norm{P_{v}\p_{\nu}U_\ell}_{L^2(\p\Si)}^2\\
        & \le -\half\norm{\p_t u_\ell}_{L^2(\p\Si)}^2,\nonumber
    \end{align}
    where the last step follows since $ P_{\cdot}$ is an orthogonal projection on $N_\eta$, and hence 
    \begin{align*}
        \norm{\p_t u_\ell}_{L^2(\p\Si)}^2&=\norm{(\ep + P_{v})\p_{\nu}U_\ell}_{L^2(\p\Si)}^2
        = \ep^2\norm{\p_{\nu} U_\ell}_{L^2(\p\Si)}^2 + (1 + 2\ep)\norm{P_{v}\p_{\nu}U_\ell}_{L^2(\p\Si)}^2,
    \end{align*} 
    and since $\ep\le \hf$. We can hence deduce that 
    \begin{equation}\label{est:galerkin_uniform}
        \sup_{0\le t\le T}\norm{\nabla U_\ell(t)}_{L^2(\Si)}^2 + \norm{\p_t u_\ell}_{L^2([0,T]\times\p\Si)}^2 + \ep\norm{\p_{\nu}U_\ell}_{L^2([0,T]\times\p\Si)}^2\leq 3\norm{\nabla U_\ell(0)}_{L^2(\Si)}^2
    \end{equation} 
    which also ensures that
    \begin{equation*}
        \begin{split}
            \norm{U_\ell(t)}_{L^2(\Si)} &\le C\norm{u_\ell(t)}_{L^2(\p\Si)} \le C\big(\norm{u_\ell(0)}_{L^2(\p\Si)} + T^{\half}\norm{\p_t u_\ell}_{L^2([0,T]\times\p\Si)}\big)\\
            &\le C\norm{U_\ell(0)}_{H^1(\Si)}
        \end{split}
    \end{equation*}
    for every $t\in [0,T]$, for a constant that is allowed to depend on $T$ and a lower bound $\iota_0$ on $\inj(g)$, but that is independent of $\ell$. 
    Hence we have that
    \beq\label{est:H1-Galerkin}
        \norm{ U_\ell}_{L^\infty([0,T];H^1(\Si))}^2+\norm{\p_t u_\ell}_{L^2([0,T]\times \p\Si)}^2 + \ep\norm{\p_{\nu} U_\ell}_{L^2([0,T]\times\p\Si)}^2\le C\norm{U_\ell(0)}_{H^1(\Si)}^2,
    \eeq
    again with a constant $C=C(T,\iota_0)$.  
    
    By construction, $\norm{u_\ell(0)}_{L^2(\p\Si)}\leq\norm{u_0}_{L^2(\p\Si)}$ since $u_\ell(0)$ is the $L^2$-orthogonal projection of $u_0$ onto $V_\ell$. Importantly, the analogous $\dot H^1(\Si)$-estimate $\norm{\nabla U_\ell(0)}_{L^2(\Si)}\leq \norm{\na (u_0)_g}_{L^2(\Si)}$ also holds, since the relation 
    \beqs
        \int_\Si\langle\nabla (\vph_j)_g,\nabla w_g\rangle\dd v = \int_{\p\Si}\p_{\nu}(\vph_j)_g w\,\dd s = \lambda_j\int_{\p\Si}\vph_j w \, \dd s \text{ for every } w\in H^\half(\pSi)
    \eeqs
    ensures that the $\la_j^{-\hf} (\vph_j)_g$, $j \geq 1$, are $\dot H^1(\Si)$-orthonormal and that $U_\ell(0)-a_0(u_0)$ can alternatively be obtained as the $\dot H^1(\Si)$-orthogonal projection of $(u_0)_g$ onto the span of $(\vphi_1)_g,\dots,(\vphi_\ell)_g$.
 
    Combining the resulting uniform bounds 
    \beqas
        \norm{U_\ell(0)}_{H^1(\Si)}&\leq C(\norm{U_\ell(0)}_{L^2(\pSi)}+\norm{\na U_\ell(0)}_{L^2(\Si)})\\
        &\leq C(\norm{u_0}_{L^2(\pSi)}+\norm{\na (u_0)_g}_{L^2(\Si)})\leq C\norm{(u_0)_g}_{H^1(\Si)}
    \eeqas
    on the initial maps $U_\ell(0)$ with the $\ell$-independent estimate \eqref{est:H1-Galerkin} thus allows us to conclude that the functions $U_\ell=(u_\ell)_g$ are uniformly bounded in $L^\infty([0,T];H^1(\Si))$ and, as $\eps>0$ is fixed, that the functions $\p_t u_\ell$ and $\p_{\nu_g}U_\ell$ are uniformly bounded in $L^2([0,T]\times\p\Si)$. This allows us to pass to a subsequence, still indexed by $\ell$, along which $U_\ell$ converges $\text{weak-}*$ in $L^\infty([0,T];H^1(\Si))$ to a function $U=(u)_g$ and along which $\p_t u_\ell$ as well as $\p_{\nu_g}U_\ell$ converge weakly in $L^2([0,T]\times\p\Si)$ to $\p_t u$ and $\p_{\nu_g}U$, respectively.

    This function hence has the regularity asked for in the lemma and since $u_\ell$ satisfies 
    \begin{equation}\label{eq:weak_PDE_const}
        \int_0^T \int_{\p\Si}\p_t u_\ell \cdot w \dd s_g \dd t = \int_0^T \int_{\p\Si} Q(\p_\nu (u_\ell)_g) \cdot w\dd s_g \dd t
    \end{equation}
    for all $w \in C^{\infty}([0,T];V_\ell)$, we find that $u$ satisfies \eqref{eq:weak_PDE_const} for all $w \in \bigcup_{j=1}^\infty C^{\infty}([0,T];V_j)$ and is hence a weak solution of \eqref{eq:ep_v_map} as required. 

    
    We furthermore observe that for solutions $u$ of this regularity, the computation carried out in \eqref{est:galerkin_uniform_bound} yields the claimed estimate \eqref{eq:energy-dec_ep_approx} on the decay of the energy. Finally we note that as our equation is \textit{linear} this in turn ensures the uniqueness of the solution in the considered class of functions.
\end{proof}

\end{document}
